% =====================================================================
%  BAB 2: MANAJEMEN KONTEN UTAMA
%  (Berita / Posts, Pengumuman / Announcements, Agenda / Events)
% =====================================================================

\chapter{Manajemen Konten Utama}

Modul ini mencakup tiga jenis konten informasi publik paling penting
dalam website sekolah: \textbf{Berita}, \textbf{Pengumuman}, dan \textbf{Agenda}.
Ketiganya dapat dikelola oleh pengguna berperan \textbf{Admin} atau \textbf{Editor}
(khusus Berita).

% -------------------------------------------------------------------
\section{Modul Berita (\textit{Posts})}
% -------------------------------------------------------------------

Modul Berita adalah sistem penerbitan artikel berbasis \textit{WYSIWYG editor}
(menggunakan \textbf{Quill Editor}). Setiap artikel mendukung fitur kategori,
\textit{tag}, thumbnail, SEO, penjadwalan tayang otomatis, dan penandaan sebagai
artikel unggulan.

\begin{figure}[H]
  \centering
  \includegraphics[width=0.9\textwidth]{../Img/Backend/screencapture-127-0-0-1-8000-admin-posts-2026-06-16-17_51_29.png}
  \caption{Daftar Berita di Panel Admin}
  \label{fig:posts-index}
\end{figure}

\subsection{Skenario: Menerbitkan Berita Prestasi Sekolah}

Misalkan SMA Nusantara baru saja meraih Juara 1 Olimpiade Sains Nasional.
Guru Pembimbing (berperan \textit{Editor}) menulis draf, kemudian Admin
menambahkan thumbnail dan memublikasikannya.

\begin{figure}[H]
  \centering
  \includegraphics[width=0.9\textwidth]{../Img/Backend/screencapture-127-0-0-1-8000-admin-posts-22-edit-2026-06-16-17_51_38.png}
  \caption{Form Editor Berita --- Tampilan Lengkap}
  \label{fig:posts-edit}
\end{figure}

\subsection{Panduan: Menulis Berita Baru}

\begin{enumerate}
  \item Klik menu \textbf{Berita} di \textit{sidebar} kiri, lalu klik tombol \textbf{Tulis Berita Baru}.
  \item Pada panel kiri atas, isi kolom \textbf{Judul} (contoh: ``Juara 1 OSN Bidang Matematika'').
        Kolom \textbf{Slug} akan terisi otomatis berdasarkan judul (dapat diubah manual jika diperlukan).
  \item Tuliskan isi lengkap artikel pada area \textbf{Konten} menggunakan Quill Editor.
        Editor mendukung pemformatan teks tebal, miring, daftar, tautan, dan penyisipan gambar.
  \item Isi kolom \textbf{Excerpt} --- ringkasan singkat artikel (1-3 kalimat) yang
        akan tampil di halaman daftar berita.
  \item Pada panel kanan, atur \textbf{Status} melalui \textit{dropdown}:
        \begin{itemize}
          \item \textit{Draf} --- tersimpan, tidak terlihat publik
          \item \textit{Review} --- menunggu persetujuan Admin
          \item \textit{Terjadwal} --- tayang otomatis pada waktu yang ditentukan
          \item \textit{Terbit} --- langsung terlihat publik
          \item \textit{Arsip} --- disembunyikan dari publik tanpa dihapus
        \end{itemize}
  \item Jika memilih \textit{Terjadwal}, isi \textbf{Jadwal Tayang} (format: Tanggal dan Jam).
  \item Centang \textbf{Tandai sebagai unggulan} jika artikel ini harus tampil di
        \textit{carousel} utama halaman depan.
  \item Pilih \textbf{Kategori} dari \textit{dropdown} yang tersedia.
  \item Tambahkan \textbf{Tag}: ketikkan nama \textit{tag} pada kolom input lalu klik tombol \textbf{Tambah},
        atau klik salah satu \textit{tag} yang sudah ada untuk memilihnya. Tag yang dipilih akan muncul
        sebagai \textit{badge} biru yang bisa dihapus dengan klik tanda \textbf{×}.
  \item Pilih gambar \textbf{Thumbnail}: klik tombol \textbf{Pilih Media} untuk membuka
        Pustaka Media. Disarankan menggunakan rasio \textbf{16:9} untuk hasil terbaik.
        Klik \textbf{Hapus} untuk membatalkan pilihan thumbnail.
  \item (Opsional) Pilih \textbf{OG Image} (Open Graph Image) untuk pratinjau saat artikel
        dibagikan di media sosial. Disarankan rasio \textbf{1.91:1}.
  \item Isi panel \textbf{SEO} di bagian bawah: \textit{Meta Title}, \textit{Meta Description},
        \textit{Canonical URL}. Aktifkan \textit{No Index} jika tidak ingin artikel ini diindeks
        oleh mesin pencari.
  \item Klik salah satu tombol aksi di \textit{action bar} yang menempel di bagian bawah layar:
        \begin{itemize}
          \item \textbf{Simpan Draf} (hitam) --- menyimpan sebagai draf
          \item \textbf{Simpan Perubahan} (biru) --- menyimpan dengan status yang dipilih
          \item \textbf{Publikasikan} (biru border) --- langsung mengubah status menjadi \textit{Published}
          \item \textbf{Preview} --- membuka pratinjau artikel di tab baru
        \end{itemize}
\end{enumerate}

\begin{notebox}[Izinkan Komentar]
  Centang opsi \textbf{Izinkan komentar} jika Anda ingin mengaktifkan fitur
  komentar pembaca pada artikel ini. Fitur ini bersifat per-artikel.
\end{notebox}

\subsection{Mengelola Kategori \& Tag}

\textbf{Kategori} dan \textbf{Tag} dikelola melalui menu terpisah di \textit{sidebar}:

\begin{enumerate}
  \item Klik menu \textbf{Kategori} untuk menambah, mengedit, atau menghapus kategori.
        Setiap kategori memiliki nama dan \textit{slug} unik.
  \item Klik menu \textbf{Tag} untuk mengelola tag. Tag bisa juga ditambahkan
        langsung dari dalam form penulisan berita.
\end{enumerate}

\subsection{Penanganan Masalah: Berita}

\begin{itemize}
  \item \textbf{Artikel tidak muncul di halaman publik}: Periksa apakah
        \textbf{Status} sudah diatur ke \textit{Terbit}. Status \textit{Draf},
        \textit{Review}, dan \textit{Arsip} tidak terlihat oleh publik.
  \item \textbf{Gambar Thumbnail tidak berubah}: Coba \textit{hard refresh} halaman
        (Ctrl+Shift+R). Perubahan gambar memerlukan pembersihan \textit{cache} browser.
  \item \textbf{Tombol Publikasikan tidak berfungsi}: Pastikan kolom \textbf{Judul}
        dan \textbf{Konten} sudah diisi. Kedua kolom ini bersifat wajib (\textit{required}).
\end{itemize}

% -------------------------------------------------------------------
\section{Modul Pengumuman (\textit{Announcements})}
% -------------------------------------------------------------------

Modul Pengumuman digunakan untuk menerbitkan pemberitahuan resmi sekolah,
seperti jadwal libur, informasi ujian, atau edaran penting. Pengumuman mendukung
fitur \textit{Pin} (sematkan) dan lampiran dokumen.

\begin{figure}[H]
  \centering
  \includegraphics[width=0.9\textwidth]{../Img/Backend/screencapture-127-0-0-1-8000-admin-announcements-2026-06-16-17_51_50.png}
  \caption{Daftar Pengumuman di Panel Admin}
  \label{fig:announcements-index}
\end{figure}

\begin{figure}[H]
  \centering
  \includegraphics[width=0.9\textwidth]{../Img/Backend/screencapture-127-0-0-1-8000-admin-announcements-1-edit-2026-06-16-17_52_05.png}
  \caption{Form Edit Pengumuman}
  \label{fig:announcements-edit}
\end{figure}

\subsection{Skenario: Menerbitkan Pengumuman Libur Nasional}

Kepala Sekolah meminta Admin TU untuk segera menerbitkan pengumuman bahwa sekolah
libur pada tanggal 17 Agustus. Admin membuat pengumuman baru, menyematkannya
(\textit{Pin}) agar muncul paling atas, dan melampirkan surat edaran PDF.

\subsection{Panduan: Membuat Pengumuman}

\begin{enumerate}
  \item Buka menu \textbf{Pengumuman} lalu klik tombol \textbf{Buat Pengumuman}.
  \item Isi kolom \textbf{Judul} (contoh: ``Libur HUT RI ke-80'').
  \item Pilih \textbf{Kategori} dari \textit{dropdown} (misal: ``Akademik'', ``Umum'').
        Pilih kosong jika tidak berkategori.
  \item Tentukan rentang aktif pengumuman:
        \begin{itemize}
          \item \textbf{Mulai} --- tanggal pengumuman mulai aktif (format: YYYY-MM-DD)
          \item \textbf{Selesai} --- tanggal pengumuman berakhir
        \end{itemize}
  \item Atur \textbf{Status}:
        \begin{itemize}
          \item \textit{Draft} --- belum terlihat publik
          \item \textit{Published} --- langsung terlihat publik
        \end{itemize}
  \item Centang \textbf{Pin Pengumuman} jika pengumuman ini harus muncul paling
        atas dalam daftar dan diberikan sorotan khusus.
  \item Tulis isi pengumuman di editor \textbf{Konten} (Quill Editor, tinggi 280px).
  \item (Opsional) Unggah berkas di kolom \textbf{Lampiran (PDF/DOC)}.
        Jika pengumuman sudah pernah memiliki lampiran dan Anda ingin menggantinya,
        centang kotak \textbf{Hapus lampiran} sebelum mengunggah file baru.
  \item Klik tombol hitam \textbf{Simpan Pengumuman}.
\end{enumerate}

\subsection{Penanganan Masalah: Pengumuman}

\begin{itemize}
  \item \textbf{Pengumuman hilang dari website sebelum tanggal selesai}: Pastikan
        \textbf{Status} masih \textit{Published}. Jika Anda atau orang lain tidak sengaja
        mengubah status ke \textit{Draft}, pengumuman akan tersembunyi.
  \item \textbf{File lampiran gagal diunggah}: Pastikan file berformat PDF atau DOC/DOCX
        dan ukurannya tidak melebihi batas unggah server (biasanya 2 MB untuk \textit{shared hosting}).
\end{itemize}

% -------------------------------------------------------------------
\section{Modul Agenda (\textit{Events})}
% -------------------------------------------------------------------

Modul Agenda mencatat jadwal kegiatan sekolah (seminar, lomba, ujian, perlombaan, dll.)
beserta lokasi, durasi, dan gambar acara.

\begin{figure}[H]
  \centering
  \includegraphics[width=0.9\textwidth]{../Img/Backend/screencapture-127-0-0-1-8000-admin-events-2026-06-16-17_52_13.png}
  \caption{Daftar Agenda di Panel Admin}
  \label{fig:events-index}
\end{figure}

\begin{figure}[H]
  \centering
  \includegraphics[width=0.9\textwidth]{../Img/Backend/screencapture-127-0-0-1-8000-admin-events-1-edit-2026-06-16-17_52_26.png}
  \caption{Form Edit Agenda --- Detail Waktu dan Lokasi}
  \label{fig:events-edit}
\end{figure}

\subsection{Panduan: Menambahkan Agenda}

\begin{enumerate}
  \item Masuk ke menu \textbf{Agenda} lalu klik tombol tambah.
  \item Isi \textbf{Judul} agenda (contoh: ``Ujian Tengah Semester Ganjil 2025'').
  \item Isi \textbf{Lokasi} tempat pelaksanaan (contoh: ``Aula SMA Nusantara'').
  \item Tentukan waktu mulai pada kolom \textbf{Mulai} (format: tanggal dan jam, contoh: 2025-10-14T08:00).
  \item Tentukan waktu selesai pada kolom \textbf{Selesai} (format sama).
  \item Tulis deskripsi lengkap agenda pada editor \textbf{Deskripsi}.
  \item Atur \textbf{Status}: \textit{Draft} atau \textit{Published}.
  \item Centang \textbf{Agenda Unggulan} jika kegiatan ini berskala besar dan
        ingin ditampilkan secara mencolok di halaman agenda publik.
  \item Unggah gambar/poster acara pada kolom \textbf{Gambar Agenda} melalui komponen
        \textit{image-upload}.
  \item Klik tombol hitam \textbf{Simpan}.
\end{enumerate}

\subsection{Penanganan Masalah: Agenda}

\begin{itemize}
  \item \textbf{Agenda tidak muncul di halaman depan}: Halaman depan biasanya menampilkan
        agenda yang memiliki tanggal \textit{Mulai} di masa mendatang. Agenda yang sudah
        lewat (\textit{expired}) tidak akan ditampilkan di \textit{upcoming events}.
  \item \textbf{Format waktu salah}: Kolom \textit{Mulai} dan \textit{Selesai} menggunakan
        format \texttt{YYYY-MM-DDTHH:mm}. Pastikan browser Anda menampilkan
        \textit{datetime-local picker} dengan benar.
\end{itemize}
